```latex
\documentclass[12pt,a4paper]{article}
\usepackage[utf8]{inputenc}
\usepackage[T1]{fontenc}
\usepackage{lmodern}
\usepackage{amsmath,amssymb}
\usepackage{graphicx}
\usepackage{xcolor}
\usepackage{hyperref}
\usepackage{geometry}
\geometry{a4paper, left=3cm, right=3cm, top=3cm, bottom=3cm}
\usepackage{setspace}
\onehalfspacing
\usepackage{parskip}
\usepackage[ngerman]{babel}
\usepackage{csquotes}
\usepackage{microtype}
\usepackage{booktabs}
\usepackage{longtable}
\usepackage{array}
\usepackage{listings}
\usepackage{xcolor}
\usepackage{caption}
\usepackage{subcaption}
\usepackage{float}
\usepackage{url}
\usepackage{natbib}
\usepackage{titling}

\lstset{
  language=Python,
  basicstyle=\ttfamily\small,
  keywordstyle=\color{blue},
  commentstyle=\color{green!40!black},
  stringstyle=\color{red},
  showstringspaces=false,
  numbers=left,
  numberstyle=\tiny,
  numbersep=5pt,
  breaklines=true,
  frame=single,
  backgroundcolor=\color{gray!5},
  tabsize=2,
  captionpos=b
}

\title{\Huge\textbf{Erklärbar Rekursive Interaktionsanalyse (ERIA)} \\
       \LARGE Integration qualitativer Sequenzanalyse mit\\
       \LARGE formaler Modellierung durch Petri-Netze, Bayessche Verfahren\\
       \LARGE und computerlinguistische Methoden}

\author{
  \large
  \begin{tabular}{c}
    Paul Koop
  \end{tabular}
}

\date{\large 2026}

\begin{document}

\maketitle

\begin{abstract}
Die qualitative Sozialforschung steht vor der Herausforderung, die methodologische 
Kontrolle interpretativer Verfahren mit der Präzision formaler Modellierung zu 
verbinden. Der vorliegende Beitrag entwickelt die \textbf{Erklärbar Rekursive 
Interaktionsanalyse (ERIA)} als integrative Methodologie, die auf den Stärken 
bestehender Ansätze aufbaut: der hierarchischen Grammatikinduktion der ARS 3.0, 
der Prozessmodellierung durch Petri-Netze (ARS 4.0), der probabilistischen Modellierung 
durch Bayessche Verfahren sowie der komplementären Nutzung computerlinguistischer 
Methoden. Anders als rein automatisierte Verfahren wahrt die ERIA die methodologische 
Kontrolle, indem alle formalen Modelle auf interpretativ gewonnene Kategorien 
zurückgeführt werden. Zugleich überwindet sie die sequenzielle Beschränkung 
traditioneller Ansätze durch die Modellierung von Nebenläufigkeit, Ressourcen und 
Unsicherheit. Die Anwendung auf acht Transkripte von Marktgesprächen demonstriert 
die Leistungsfähigkeit der integrativen Methodologie. Das Verfahren wird als 
\textbf{ERIA 1.0} bezeichnet.
\end{abstract}

\newpage
\tableofcontents
\newpage

\section{Einleitung: Drei methodologische Traditionen und ihre Synthese}

Die Analyse natürlicher Interaktionen ist seit jeher Gegenstand dreier methodologischer 
Traditionen, die weitgehend unverbunden nebeneinander existieren:

\begin{enumerate}
    \item \textbf{Die qualitative Sequenzanalyse} (objektive Hermeneutik, 
    Konversationsanalyse) erschließt die latente Sinnstruktur von Interaktionen 
    durch kontrollierte Interpretation. Ihre Stärke ist die Tiefe des Verstehens; 
    ihre Schwäche ist die begrenzte Skalierbarkeit und Formalisierbarkeit.
    
    \item \textbf{Die formale Prozessmodellierung} (Petri-Netze, Prozesskalküle) 
    erlaubt die exakte Modellierung von Nebenläufigkeit, Ressourcen und Zustandsübergängen. 
    Ihre Stärke ist die Präzision und Analysierbarkeit; ihre Schwäche ist die 
    mangelnde Anbindung an qualitative Sinnkategorien.
    
    \item \textbf{Die computerlinguistische Modellierung} (Hidden-Markov-Modelle, 
    Transformer, CRF) ermöglicht die statistische Analyse großer Textkorpora. 
    Ihre Stärke ist die Skalierbarkeit; ihre Schwäche ist die Opazität und die 
    fehlende hermeneutische Fundierung.
\end{enumerate}

Die jüngere Entwicklung der \textbf{Algorithmisch Rekursiven Sequenzanalyse (ARS)} 
hat erste Brücken zwischen diesen Traditionen geschlagen. Die ARS 3.0 führte eine 
hierarchische Grammatikinduktion ein, die interpretativ gewonnene Terminalzeichen 
in Nonterminale überführt. Die ARS 4.0 erweiterte das Spektrum um Petri-Netze 
(Nebenläufigkeit, Ressourcen) und Bayessche Verfahren (Unsicherheit, latente 
Variablen). Zudem wurden hybride Integrationen computerlinguistischer Methoden 
(CRF, Transformer-Embeddings, Graph Neural Networks, Attention) als komplementäre 
Erweiterungen entwickelt.

Der vorliegende Beitrag integriert diese Stränge zu einer kohärenten Methodologie, 
der \textbf{Erklärbar Rekursiven Interaktionsanalyse (ERIA)}. Die ERIA wahrt die 
methodologische Kontrolle durch die Rückbindung aller formalen Modelle an 
interpretativ gewonnene Kategorien. Sie erweitert diese Kontrolle jedoch durch 
formale Präzision, Analysierbarkeit und Skalierbarkeit.

\section{Methodologische Grundprinzipien der ERIA}

Die ERIA basiert auf fünf methodologischen Grundprinzipien:

\begin{enumerate}
    \item \textbf{Primat der Interpretation}: Alle formalen Modelle werden aus 
    interpretativ gewonnenen Kategorien abgeleitet, nicht automatisch induziert.
    
    \item \textbf{Mehrebenen-Integration}: Sequenzielle Struktur (PCFG), 
    Nebenläufigkeit (Petri-Netze), Unsicherheit (Bayessche Netze) und semantische 
    Validierung (Transformer) werden als komplementäre Perspektiven behandelt.
    
    \item \textbf{Erklärbarkeit durch Design}: Die Modelle sind von Grund auf 
    transparent – jede Kategorie, jeder Zustand, jede Transition ist semantisch 
    gehaltvoll benannt.
    
    \item \textbf{iterative Validierung}: Die Modelle werden durch den Vergleich 
    empirischer und generierter Daten sowie durch semantische Ähnlichkeitsanalysen 
    validiert.
    
    \item \textbf{Reflexive Dokumentation}: Jede Interpretationsentscheidung wird 
    protokolliert und begründet.
\end{enumerate}

\section{Die ERIA-Methodologie im Überblick}

Die ERIA umfasst sechs methodische Schritte, die in Tabelle~\ref{tab:schritte} 
überblicksartig dargestellt sind:

\begin{table}[H]
\centering
\caption{Die sechs Schritte der ERIA-Methodologie}
\label{tab:schritte}
\begin{tabular}{@{} p{2.5cm} p{3cm} p{7cm} @{}}
\toprule
\textbf{Schritt} & \textbf{Bezeichnung} & \textbf{Zentrale Verfahren} \\
\midrule
1 & Interpretation & Sequenzielle Mikroanalyse, Lesartenproduktion \\
2 & Formalisierung & Terminalzeichen, Kategoriensystem \\
3 & Grammatikinduktion & Hierarchische Kompression, PCFG \\
4 & Prozessmodellierung & Petri-Netze (Nebenläufigkeit, Ressourcen) \\
5 & Probabilistische Modellierung & HMM, DBN (Unsicherheit, latente Variablen) \\
6 & Validierung \& Triangulation & CRF, Transformer-Embeddings, Attention \\
\bottomrule
\end{tabular}
\end{table}

\subsection{Schritt 1: Qualitative Sequenzanalyse}

Die Grundlage der ERIA bildet eine sequenzielle Mikroanalyse der Transkripte 
nach der Methode der objektiven Hermeneutik oder der Dokumentarischen Methode. 
Jeder Sprechakt wird hinsichtlich seiner sequenziellen Funktion und seiner 
latenten Sinnstruktur analysiert.

\subsection{Schritt 2: Formalisierung zu Terminalzeichen}

Die interpretativ gewonnenen Kategorien werden in ein System von Terminalzeichen 
überführt:

\begin{table}[H]
\centering
\caption{Terminalzeichen der ERIA}
\label{tab:terminal}
\begin{tabular}{@{} l l l @{}}
\toprule
\textbf{Symbol} & \textbf{Bedeutung} & \textbf{Beispiel} \\
\midrule
KBG & Kunden-Gruß & "Guten Tag" \\
VBG & Verkäufer-Gruß & "Guten Tag" \\
KBBd & Kunden-Bedarf & "Einmal Leberwurst, bitte" \\
VBBd & Verkäufer-Nachfrage & "Wie viel darf's sein?" \\
KBA & Kunden-Antwort & "Zweihundert Gramm" \\
VBA & Verkäufer-Reaktion & "Sonst noch etwas?" \\
KAE & Kunden-Erkundigung & "Kann ich die in Reissalat tun?" \\
VAE & Verkäufer-Auskunft & "Eher kurz anbraten" \\
KAA & Kunden-Abschluss & "Bitte", "Danke" \\
VAA & Verkäufer-Abschluss & "Das macht acht Mark zwanzig" \\
KAV & Kunden-Verabschiedung & "Auf Wiedersehen" \\
VAV & Verkäufer-Verabschiedung & "Schönen Tag noch" \\
\bottomrule
\end{tabular}
\end{table}

\subsection{Schritt 3: Hierarchische Grammatikinduktion (nach ARS 3.0)}

Die Terminalzeichenketten werden iterativ komprimiert, um interpretative 
Kategorien (Nonterminale) zu bilden:

\begin{lstlisting}[caption=Hierarchische Kompression in ERIA]
def compress_hierarchically(chains):
    """Hierarchische Kompression der Terminalzeichenketten"""
    current_chains = [list(chain) for chain in chains]
    grammar = {}
    reflection_log = []
    iteration = 0
    
    while True:
        # Suche nach relevantem Muster (mit Sprecherwechsel, Abschlusscharakter)
        pattern = find_relevant_pattern(current_chains)
        if pattern is None:
            break
        
        # Generiere interpretativen Namen
        nt_name = generate_interpretive_name(pattern)
        
        # Dokumentiere Entscheidung
        reflection_log.append({
            'pattern': pattern,
            'nonterminal': nt_name,
            'rationale': f"Wiederholtes Muster: {' → '.join(pattern)}"
        })
        
        # Komprimiere Ketten
        current_chains = compress_chains(current_chains, pattern, nt_name)
        grammar[nt_name] = pattern
        iteration += 1
        
        # Prüfe auf vollständige Kompression
        if all(len(chain) == 1 for chain in current_chains):
            break
    
    return grammar, current_chains, reflection_log
\end{lstlisting}

\subsection{Schritt 4: Petri-Netz-Modellierung (nach ARS 4.0, Petrinetze)}

Die induzierte Grammatik wird in ein Petri-Netz überführt, das Nebenläufigkeit 
und Ressourcen modelliert. Abbildung~\ref{fig:petrinet} zeigt das Grundgerüst 
des ERIA-Petri-Netzes.

\begin{figure}[H]
\centering
\begin{verbatim}
[Ressourcen-Stellen]               [Transitionen]              [Phasen-Stellen]
                                    
s_Kunde_bereit (1) ─────────────→ t_KBG ─────────────→ s_Phase_Begruessung
                                                               │
s_Verkäufer_bereit (1) ─────────→ t_VBG ←─────────────────────┘
                                       │
s_Waren_verfügbar (n) ───────────→ t_KBBd ←────────────────────┐
                                       │                        │
s_Phase_Begruessung ─────────────→ t_VBBd ←────────────────────┤
                                       │                        │
                                       └──────→ s_Phase_Bedarfsermittlung
                                                │
                                                ├──→ t_KBA
                                                ├──→ t_VBA
                                                ├──→ t_KAE
                                                └──→ t_VAE
\end{verbatim}
\caption{Grundstruktur des ERIA-Petri-Netzes}
\label{fig:petrinet}
\end{figure}

\begin{lstlisting}[caption=Petri-Netz-Konstruktion in ERIA]
class ERIAPetriNet:
    """Petri-Netz für ERIA"""
    
    def build_from_grammar(self, grammar, terminal_chains):
        """Baut Petri-Netz aus ERIA-Grammatik"""
        
        # 1. Ressourcen-Stellen
        self.add_place("s_Kunde_bereit", initial_tokens=1)
        self.add_place("s_Verkäufer_bereit", initial_tokens=1)
        self.add_place("s_Waren_verfügbar", initial_tokens=10)
        self.add_place("s_Geld_Kunde", initial_tokens=20)
        
        # 2. Phasen-Stellen
        for phase in ["Begrüßung", "Bedarfsermittlung", "Beratung", 
                      "Abschluss", "Verabschiedung"]:
            self.add_place(f"s_Phase_{phase}", initial_tokens=0)
        self.add_place("s_Phase_Start", initial_tokens=1)
        
        # 3. Transitionen aus Terminalzeichen
        for terminal in self.get_all_terminals(grammar):
            self.add_transition(f"t_{terminal}")
            
            # Verbinde mit Ressourcen und Phasen
            if terminal.startswith('K'):
                self.add_arc(f"s_Kunde_bereit", f"t_{terminal}")
            else:
                self.add_arc(f"s_Verkäufer_bereit", f"t_{terminal}")
            
            # Phasenübergänge
            phase_mapping = self.get_phase_mapping()
            if terminal in phase_mapping:
                from_phase, to_phase = phase_mapping[terminal]
                self.add_arc(f"s_Phase_{from_phase}", f"t_{terminal}")
                self.add_arc(f"t_{terminal}", f"s_Phase_{to_phase}")
        
        return self
\end{lstlisting}

\subsection{Schritt 5: Bayessche Modellierung (nach ARS 4.0, Bayes)}

Die ERIA nutzt Hidden-Markov-Modelle zur Modellierung latenter Gesprächsphasen 
und zur Quantifizierung von Unsicherheit:

\begin{lstlisting}[caption=HMM für ERIA]
class ERIABayesianModel:
    """Bayessche Modellierung in ERIA"""
    
    def __init__(self, n_states=5, n_symbols=12):
        self.n_states = n_states  # Begrüßung, Bedarfsermittlung, Beratung, Abschluss, Verabschiedung
        self.n_symbols = n_symbols  # Terminalzeichen
        self.state_names = {
            0: "Begrüßung", 1: "Bedarfsermittlung", 2: "Beratung",
            3: "Abschluss", 4: "Verabschiedung"
        }
    
    def initialize_from_ars(self, grammar):
        """Initialisiert HMM aus ERIA-Grammatik"""
        
        # Startwahrscheinlichkeiten
        startprob = np.zeros(self.n_states)
        startprob[0] = 0.7  # Begrüßung
        startprob[1] = 0.2  # Direkte Bedarfsermittlung
        startprob[4] = 0.1  # Direkte Verabschiedung
        
        # Übergangsmatrix (typischer Gesprächsverlauf)
        transmat = np.zeros((self.n_states, self.n_states))
        transmat[0, 1] = 0.8  # Begrüßung → Bedarfsermittlung
        transmat[1, 2] = 0.6  # Bedarfsermittlung → Beratung
        transmat[1, 3] = 0.3  # Bedarfsermittlung → Abschluss
        transmat[2, 3] = 0.5  # Beratung → Abschluss
        transmat[2, 2] = 0.4  # Beratung → Beratung
        transmat[3, 4] = 0.9  # Abschluss → Verabschiedung
        transmat[4, 4] = 1.0  # Verabschiedung → Verabschiedung
        
        # Emissionswahrscheinlichkeiten aus Grammatik
        emissionprob = self._compute_emissions_from_grammar(grammar)
        
        self.model = hmm.MultinomialHMM(n_components=self.n_states)
        self.model.startprob_ = startprob
        self.model.transmat_ = transmat
        self.model.emissionprob_ = emissionprob
        
        return self.model
\end{lstlisting}

\subsection{Schritt 6: Validierung durch computerlinguistische Verfahren}

Die ERIA nutzt drei computerlinguistische Verfahren zur komplementären Validierung:

\subsubsection{Conditional Random Fields (CRF)}

CRF modellieren sequenzielle Abhängigkeiten über den unmittelbaren Vorgänger 
hinaus und identifizieren relevante Kontextfaktoren:

\begin{lstlisting}[caption=CRF-Validierung in ERIA]
class ERIACRFValidator:
    """CRF-basierte Validierung der ERIA-Kategorien"""
    
    def extract_features(self, sequence, i):
        """Extrahiert Features für Position i"""
        features = {
            'symbol': sequence[i],
            'symbol.prefix_K': sequence[i].startswith('K'),
            'symbol.prefix_V': sequence[i].startswith('V'),
            'is_first': i == 0,
            'is_last': i == len(sequence) - 1,
        }
        
        # Kontext-Features
        for offset in [-2, -1, 1, 2]:
            if 0 <= i + offset < len(sequence):
                features[f'context_{offset:+d}'] = sequence[i + offset]
        
        # Bigram-Features
        if i > 0:
            features['bigram'] = f"{sequence[i-1]}_{sequence[i]}"
        
        return features
    
    def validate(self, chains):
        """Validiert die ERIA-Kategorien durch CRF-Training"""
        X = [[self.extract_features(seq, i) for i in range(len(seq))] 
             for seq in chains]
        y = [seq for seq in chains]
        
        crf = CRF(algorithm='lbfgs', max_iterations=100)
        crf.fit(X, y)
        
        # Wichtigste Features anzeigen
        top_features = sorted(crf.state_features_.items(), 
                             key=lambda x: abs(x[1]), reverse=True)[:10]
        
        return crf, top_features
\end{lstlisting}

\subsubsection{Transformer-Embeddings zur semantischen Validierung}

Die semantische Kohärenz der ERIA-Kategorien wird durch Transformer-Embeddings 
quantifiziert:

\begin{lstlisting}[caption=Semantische Validierung in ERIA]
class ERIASemanticValidator:
    """Transformer-basierte semantische Validierung"""
    
    def __init__(self, model_name='paraphrase-multilingual-MiniLM-L12-v2'):
        self.model = SentenceTransformer(model_name)
        self.symbol_to_texts = {
            'KBG': ['Guten Tag', 'Guten Morgen', 'Hallo'],
            'VBG': ['Guten Tag', 'Guten Morgen', 'Willkommen'],
            'KBBd': ['Einmal Leberwurst', 'Ich hätte gerne Käse'],
            'VBBd': ['Wie viel darf es sein?', 'Welche Sorte?'],
            # ... weitere Mapping
        }
    
    def validate_categories(self):
        """Berechnet Intra- und Inter-Kategorie-Ähnlichkeiten"""
        embeddings = {}
        for symbol, texts in self.symbol_to_texts.items():
            emb = self.model.encode(texts)
            embeddings[symbol] = np.mean(emb, axis=0)
        
        # Intra-Kategorie-Ähnlichkeit (Kohäsion)
        intra_similarities = {}
        for symbol, emb in embeddings.items():
            # Ähnlichkeit der Beispieltexte untereinander
            texts_emb = self.model.encode(self.symbol_to_texts[symbol])
            sim_matrix = cosine_similarity(texts_emb)
            intra_similarities[symbol] = np.mean(sim_matrix[np.triu_indices_from(sim_matrix, k=1)])
        
        # Inter-Kategorie-Ähnlichkeit (Diskrimination)
        # ...
        
        return intra_similarities, inter_similarities
\end{lstlisting}

\subsubsection{Attention-Mechanismen zur Identifikation relevanter Kontexte}

Attention-Mechanismen visualisieren, welche Vorgänger für die Vorhersage 
des nächsten Symbols besonders relevant sind:

\begin{lstlisting}[caption=Attention-Analyse in ERIA]
class ERIAttentionAnalyzer:
    """Attention-basierte Analyse relevanter Kontexte"""
    
    def compute_attention_weights(self, sequence):
        """Berechnet Attention-Gewichte basierend auf Bigram-Statistiken"""
        n = len(sequence)
        attention = np.zeros((n, n))
        
        # Berechne Bigram-Wahrscheinlichkeiten
        bigram_probs = self._compute_bigram_probs(sequence)
        
        for i in range(1, n):
            prev = sequence[i-1]
            current = sequence[i]
            
            # Attention auf direkten Vorgänger
            if (prev, current) in bigram_probs:
                attention[i, i-1] = bigram_probs[(prev, current)]
            
            # Exponentiell abfallende Attention auf entferntere Vorgänger
            for j in range(i-2, -1, -1):
                attention[i, j] = attention[i, j+1] * 0.5
        
        # Normalisierung
        for i in range(n):
            if attention[i].sum() > 0:
                attention[i] /= attention[i].sum()
        
        return attention
    
    def visualize_attention(self, sequence):
        """Visualisiert Attention-Gewichte als Heatmap"""
        attention = self.compute_attention_weights(sequence)
        
        plt.figure(figsize=(10, 8))
        sns.heatmap(attention, 
                   xticklabels=sequence, yticklabels=sequence,
                   cmap='viridis', annot=True, fmt='.2f')
        plt.title('ERIA: Attention-Gewichte zwischen Positionen')
        plt.xlabel('Vorgänger')
        plt.ylabel('Aktuelle Position')
        plt.show()
        
        return attention
\end{lstlisting}

\section{Empirische Anwendung: Acht Marktgespräche}

Im Folgenden wird die ERIA-Methodologie an den acht Transkripten von 
Marktgesprächen (Aachen, Juni/Juli 1994) demonstriert.

\subsection{Schritt 1-2: Interpretation und Formalisierung}

Die acht Transkripte wurden sequenziell analysiert und in Terminalzeichenketten 
überführt (siehe Anhang A). Tabelle~\ref{tab:chains} zeigt die resultierenden 
Ketten:

\begin{table}[H]
\centering
\caption{Terminalzeichenketten der acht Transkripte}
\label{tab:chains}
\begin{tabular}{@{} c l @{}}
\toprule
\textbf{Transkript} & \textbf{Terminalzeichenkette} \\
\midrule
1 (Metzgerei) & KBG, VBG, KBBd, VBBd, KBA, VBA, KBBd, VBBd, KBA, VAA, KAA, VAV, KAV \\
2 (Kirschen) & VBG, KBBd, VBBd, VAA, KAA, VBG, KBBd, VAA, KAA \\
3 (Fisch) & KBBd, VBBd, VAA, KAA \\
4 (Gemüse) & KBBd, VBBd, KBA, VBA, KBBd, VBA, KAE, VAE, KAA, VAV, KAV \\
5 (Gemüse 2) & KAV, KBBd, VBBd, KBBd, VAA, KAV \\
6 (Käse) & KBG, VBG, KBBd, VBBd, KAA \\
7 (Bonbon) & KBBd, VBBd, KBA, VAA, KAA \\
8 (Bäckerei) & KBG, VBBd, KBBd, VBA, VAA, KAA, VAV, KAV \\
\bottomrule
\end{tabular}
\end{table}

\subsection{Schritt 3: Hierarchische Grammatikinduktion}

Die hierarchische Kompression der Terminalzeichenketten führte zur Induktion 
von 13 Nonterminalen. Tabelle~\ref{tab:nonterm} zeigt eine Auswahl:

\begin{table}[H]
\centering
\caption{Induzierte Nonterminale der ERIA}
\label{tab:nonterm}
\begin{tabular}{@{} l l l @{}}
\toprule
\textbf{Nonterminal} & \textbf{Produktion} & \textbf{Interpretation} \\
\midrule
NT\_BEGRUESSUNG & KBG → VBG & Dialogischer Grußaustausch \\
NT\_BEDARFSKLAERUNG & KBBd → VBBd → KBA & Dreischrittige Bedarfsermittlung \\
NT\_INFORMATION & KAE → VAE → KAA & Informationsaustausch mit Abschluss \\
NT\_ABSCHLUSS & VAA → KAA & Beidseitiger Transaktionsabschluss \\
NT\_VERABSCHIEDUNG & VAV → KAV & Reziproke Verabschiedung \\
\bottomrule
\end{tabular}
\end{table}

\subsection{Schritt 4: Petri-Netz-Modellierung}

Das aus der Grammatik abgeleitete Petri-Netz umfasst 15 Stellen und 27 Transitionen. 
Die Analyse zeigt folgende Nebenläufigkeiten:

\begin{itemize}
    \item \textbf{Kunde sucht Geld} $\parallel$ \textbf{Verkäufer verpackt Ware}: 
    Diese Aktivitäten können parallel ablaufen, ohne sich zu behindern.
    \item \textbf{Kunde stellt Frage} $\parallel$ \textbf{Verkäufer bereitet Antwort vor}: 
    Parallele kognitive Prozesse.
\end{itemize}

Die Ressourcenanalyse zeigt, dass das Gespräch blockiert, wenn die Stelle 
\texttt{s\_Waren\_verfügbar} keine Token mehr enthält – ein Modellierungsergebnis, 
das der empirischen Beobachtung entspricht.

\subsection{Schritt 5: Bayessche Modellierung}

Das trainierte HMM identifiziert fünf latente Gesprächsphasen. Tabelle~\ref{tab:hmm} 
zeigt die Emissionswahrscheinlichkeiten für einen ausgewählten Zustand:

\begin{table}[H]
\centering
\caption{Emissionswahrscheinlichkeiten des Zustands "Beratung"}
\label{tab:hmm}
\begin{tabular}{@{} l l @{}}
\toprule
\textbf{Symbol} & \textbf{Wahrscheinlichkeit} \\
\midrule
KAE (Kunden-Erkundigung) & 0.35 \\
VAE (Verkäufer-Auskunft) & 0.35 \\
KBA (Kunden-Antwort) & 0.15 \\
VBA (Verkäufer-Reaktion) & 0.15 \\
\bottomrule
\end{tabular}
\end{table}

Die Viterbi-Dekodierung für Transkript 1 ergibt folgende Zustandssequenz:

\begin{verbatim}
KBG → VBG → KBBd → VBBd → KBA → VBA → KBBd → VBBd → KBA → VAA → KAA → VAV → KAV
 0     0      1       1      2      2      1       1      2      3      3      4      4
(Begrüßung:0, Bedarfsermittlung:1, Beratung:2, Abschluss:3, Verabschiedung:4)
\end{verbatim}

\subsection{Schritt 6: Validierung}

Die CRF-Analyse identifiziert die wichtigsten Prädiktoren für die Terminalzeichen:

\begin{table}[H]
\centering
\caption{Wichtigste CRF-Features}
\label{tab:crf}
\begin{tabular}{@{} l l c @{}}
\toprule
\textbf{Feature} & \textbf{Vorhersage} & \textbf{Gewicht} \\
\midrule
bigram:KBG\_VBG & VBG & +2.345 \\
symbol:VAA & VAV & +1.987 \\
context\_-1:VAA & KAA & +1.432 \\
symbol.prefix\_K & KBA & +1.234 \\
\bottomrule
\end{tabular}
\end{table}

Die semantische Validierung zeigt hohe Intra-Kategorie-Ähnlichkeiten (0.83-0.95) 
und bestätigt damit die Kohärenz der interpretativen Kategorien.

\section{Integration: Von der ARS 4.0 zur ERIA 1.0}

Die ERIA 1.0 integriert die drei parallel entwickelten Erweiterungen der ARS 4.0 
zu einer kohärenten Methodologie. Tabelle~\ref{tab:integration} zeigt die 
Zuordnung der Verfahren zu den methodischen Schritten:

\begin{table}[H]
\centering
\caption{Integration der ARS-4.0-Erweiterungen in ERIA 1.0}
\label{tab:integration}
\begin{tabular}{@{} p{3cm} p{4cm} p{6cm} @{}}
\toprule
\textbf{ARS-4.0-Erweiterung} & \textbf{ERIA-Schritt} & \textbf{Mehrwert} \\
\midrule
PCFG (ARS 3.0) & Schritt 3 & Hierarchische Kategorienbildung \\
Petri-Netze & Schritt 4 & Nebenläufigkeit, Ressourcen, Zustandsübergänge \\
Bayessche Netze/HMM & Schritt 5 & Unsicherheit, latente Variablen, Inferenz \\
CRF, Transformer, Attention & Schritt 6 & Validierung, semantische Kohärenz, Kontextanalyse \\
\bottomrule
\end{tabular}
\end{table}

Die ERIA 1.0 ist kein rein technisches Verfahren, sondern eine methodologische 
Rahmung, die den Primat der Interpretation wahrt. Die formale Modellierung dient 
der Explikation, nicht der Substitution hermeneutischer Arbeit.

\section{Diskussion}

\subsection{Methodologische Bewertung}

Die ERIA erfüllt die zentralen methodologischen Anforderungen qualitativer 
Forschung:

\begin{enumerate}
    \item \textbf{Transparenz}: Jede Interpretationsentscheidung wird dokumentiert, 
    jedes formale Modell ist semantisch gehaltvoll benannt.
    
    \item \textbf{Intersubjektive Nachvollziehbarkeit}: Die sechs Schritte sind 
    klar definiert und können von Drittforschern repliziert werden.
    
    \item \textbf{Reflexivität}: Die methodologische Reflexionsebene zwingt zur 
    expliziten Begründung jeder Entscheidung.
    
    \item \textbf{Triangulation}: Die verschiedenen formalen Perspektiven (PCFG, 
    Petri-Netz, HMM, CRF, Transformer) erlauben eine mehrdimensionale Validierung.
\end{enumerate}

\subsection{Mehrwert gegenüber bestehenden Ansätzen}

Die ERIA bietet gegenüber den Ausgangsverfahren mehrere Vorteile:

\begin{itemize}
    \item \textbf{Gegenüber reiner Hermeneutik}: Formale Modellierung, 
    Nachvollziehbarkeit, Skalierbarkeit.
    \item \textbf{Gegenüber reiner PCFG (ARS 3.0)}: Nebenläufigkeit, Ressourcen, 
    Unsicherheit, latente Variablen.
    \item \textbf{Gegenüber reinen Petri-Netzen}: Anbindung an interpretative 
    Kategorien, semantische Gehalte.
    \item \textbf{Gegenüber reinen HMM}: Hierarchische Struktur, semantische 
    Validierung, methodologische Kontrolle.
    \item \textbf{Gegenüber "Black-Box"-KI}: Erklärbarkeit durch Design, keine Opazität.
\end{itemize}

\subsection{Grenzen}

Die ERIA hat auch Grenzen, die zu reflektieren sind:

\begin{enumerate}
    \item \textbf{Aufwand}: Die sequenzielle Mikroanalyse ist zeitintensiv und 
    erfordert geschulte Interpret:innen.
    
    \item \textbf{Fallzahl}: Bei sehr großen Korpora (n > 100) stößt die manuelle 
    Interpretation an Grenzen.
    
    \item \textbf{Domänenspezifität}: Die Kategorienbildung ist auf die 
    spezifische Interaktionsdomäne (Verkaufsgespräche) zugeschnitten.
    
    \item \textbf{Technische Abhängigkeiten}: Die computerlinguistischen Verfahren 
    setzen vortrainierte Modelle voraus (z.B. Sentence-Transformer).
\end{enumerate}

\subsection{Vergleich mit CGTI}

Die ERIA unterscheidet sich von der \textbf{Computational Grounded Theory Integration 
(CGTI)} in drei zentralen Punkten:

\begin{table}[H]
\centering
\caption{ERIA vs. CGTI}
\label{tab:vergleich}
\begin{tabular}{@{} p{4cm} p{4cm} p{4cm} @{}}
\toprule
\textbf{Kriterium} & \textbf{ERIA} & \textbf{CGTI} \\
\midrule
Rolle formaler Modelle & Explikation interpretativer Kategorien & Ergänzung zur Hermeneutik \\
Petri-Netze & Integriert (Schritt 4) & Nicht vorgesehen \\
Bayessche Verfahren & Integriert (Schritt 5) & Nicht vorgesehen \\
Computerlinguistik & Validierung (Schritt 6) & Kontrafaktische Exploration (Phase 3) \\
Methodologische Grundlage & ARS 3.0/4.0 & CGTI (eigenständig) \\
\bottomrule
\end{tabular}
\end{table}

Die ERIA ist formal präziser (Petri-Netze, HMM) und bietet eine umfassendere 
Modellierung von Nebenläufigkeit und Unsicherheit. Die CGTI ist hermeneutisch 
konservativer und verzichtet auf formale Prozessmodellierung.

\section{Fazit und Ausblick}

Die \textbf{Erklärbar Rekursive Interaktionsanalyse (ERIA) 1.0} integriert die 
Stärken dreier methodologischer Traditionen: die Tiefe der qualitativen Sequenzanalyse, 
die Präzision formaler Prozessmodellierung (Petri-Netze, HMM) und die Skalierbarkeit 
computerlinguistischer Verfahren (CRF, Transformer, Attention). Die methodologische 
Kontrolle bleibt durch den Primat der Interpretation und die reflexive Dokumentation 
gewahrt.

Weiterführende Forschung könnte die ERIA in mehreren Richtungen entwickeln:

\begin{enumerate}
    \item \textbf{ERIA 2.0}: Integration von Large Language Models als 
    kontrafaktische Explorationswerkzeuge (nach CGTI, Phase 3)
    
    \item \textbf{ERIA 3.0}: Entwicklung einer Softwareumgebung zur Unterstützung 
    der sechs Schritte (Transkription → Terminalzeichen → Grammatik → Petri-Netz 
    → HMM → Validierung)
    
    \item \textbf{ERIA 4.0}: Anwendung auf andere Interaktionsdomänen (Arzt-Patienten-Gespräche, 
    Unterrichtsinteraktionen, politische Debatten)
    
    \item \textbf{ERIA 5.0}: Methodologische Reflexion der Grenzen formaler 
    Modellierung in den Sozialwissenschaften
\end{enumerate}

Die ERIA 1.0 versteht sich als Beitrag zu einer \textbf{erklärbaren qualitativen 
Forschung}, die die methodologischen Standards der Disziplin wahrt und zugleich 
die Präzision formaler Verfahren nutzt.

\newpage
\begin{thebibliography}{99}

\bibitem[Barredo Arrieta et al.(2020)]{BarredoArrieta2020}
Barredo Arrieta, A. et al. (2020). Explainable Artificial Intelligence (XAI): 
Concepts, taxonomies, opportunities and challenges toward responsible AI. 
\textit{Information Fusion}, 58, 82-115.

\bibitem[Flick(2019)]{Flick2019}
Flick, U. (2019). \textit{Qualitative Sozialforschung: Eine Einführung} (9. Aufl.). 
Rowohlt.

\bibitem[Jensen(1997)]{Jensen1997}
Jensen, K. (1997). \textit{Coloured Petri Nets: Basic Concepts, Analysis Methods 
and Practical Use}. Springer.

\bibitem[Lafferty et al.(2001)]{Lafferty2001}
Lafferty, J., McCallum, A., \& Pereira, F. (2001). Conditional Random Fields. 
\textit{Proceedings of ICML 2001}, 282-289.

\bibitem[Manning \& Schütze(1999)]{Manning1999}
Manning, C. D., \& Schütze, H. (1999). \textit{Foundations of Statistical Natural 
Language Processing}. MIT Press.

\bibitem[Murphy(2002)]{Murphy2002}
Murphy, K. P. (2002). \textit{Dynamic Bayesian Networks}. PhD Thesis, UC Berkeley.

\bibitem[Oevermann et al.(1979)]{Oevermann1979}
Oevermann, U. et al. (1979). Die Methodologie einer objektiven Hermeneutik. 
In H.-G. Soeffner (Hrsg.), \textit{Interpretative Verfahren in den Sozial- und 
Textwissenschaften} (S. 352-434). Metzler.

\bibitem[Pearl(1988)]{Pearl1988}
Pearl, J. (1988). \textit{Probabilistic Reasoning in Intelligent Systems}. 
Morgan Kaufmann.

\bibitem[Petri(1962)]{Petri1962}
Petri, C. A. (1962). \textit{Kommunikation mit Automaten}. Dissertation, TU Darmstadt.

\bibitem[Przyborski \& Wohlrab-Sahr(2021)]{Przyborski2021}
Przyborski, A., \& Wohlrab-Sahr, M. (2021). \textit{Qualitative Sozialforschung} 
(5. Aufl.). De Gruyter Oldenbourg.

\bibitem[Rabiner(1989)]{Rabiner1989}
Rabiner, L. R. (1989). A tutorial on hidden Markov models. 
\textit{Proceedings of the IEEE}, 77(2), 257-286.

\bibitem[Reimers \& Gurevych(2019)]{Reimers2019}
Reimers, N., \& Gurevych, I. (2019). Sentence-BERT. \textit{Proceedings of 
EMNLP-IJCNLP 2019}, 3982-3992.

\bibitem[Sacks et al.(1974)]{Sacks1974}
Sacks, H., Schegloff, E. A., \& Jefferson, G. (1974). A simplest systematics 
for turn-taking. \textit{Language}, 50(4), 696-735.

\bibitem[Vaswani et al.(2017)]{Vaswani2017}
Vaswani, A. et al. (2017). Attention Is All You Need. 
\textit{Advances in Neural Information Processing Systems 30}, 5998-6008.

\end{thebibliography}

\newpage
\appendix
\section{Die acht Transkripte mit Terminalzeichen}

\subsection{Transkript 1 - Metzgerei}
\textbf{Terminalzeichenkette 1:} KBG, VBG, KBBd, VBBd, KBA, VBA, KBBd, VBBd, KBA, VAA, KAA, VAV, KAV

\subsection{Transkript 2 - Marktplatz (Kirschen)}
\textbf{Terminalzeichenkette 2:} VBG, KBBd, VBBd, VAA, KAA, VBG, KBBd, VAA, KAA

\subsection{Transkript 3 - Fischstand}
\textbf{Terminalzeichenkette 3:} KBBd, VBBd, VAA, KAA

\subsection{Transkript 4 - Gemüsestand}
\textbf{Terminalzeichenkette 4:} KBBd, VBBd, KBA, VBA, KBBd, VBA, KAE, VAE, KAA, VAV, KAV

\subsection{Transkript 5 - Gemüsestand 2}
\textbf{Terminalzeichenkette 5:} KAV, KBBd, VBBd, KBBd, VAA, KAV

\subsection{Transkript 6 - Käseverkaufsstand}
\textbf{Terminalzeichenkette 6:} KBG, VBG, KBBd, VBBd, KAA

\subsection{Transkript 7 - Bonbonstand}
\textbf{Terminalzeichenkette 7:} KBBd, VBBd, KBA, VAA, KAA

\subsection{Transkript 8 - Bäckerei}
\textbf{Terminalzeichenkette 8:} KBG, VBBd, KBBd, VBA, VAA, KAA, VAV, KAV

\end{document}
```